\documentclass[a4paper,12pt]{ctexart}
\usepackage{xcolor,graphicx,geometry,array,hyperref,booktabs}
\hypersetup{colorlinks=true,linkcolor=blue!70!black}
\geometry{top=1.8cm,bottom=1.8cm,left=1.8cm,right=1.8cm}
\linespread{1.35}

\newenvironment{thinkbox}{\vspace{0.3cm}\begin{center}\begin{minipage}{0.88\textwidth}\colorbox{blue!5}{\begin{minipage}{0.94\textwidth}\vspace{0.15cm}\textbf{【想一想】}}\vspace{0.15cm}\end{minipage}\end{minipage}\end{center}\vspace{0.3cm}}{}
\newenvironment{funfact}{\vspace{0.2cm}\begin{center}\begin{minipage}{0.88\textwidth}\fbox{\begin{minipage}{0.92\textwidth}\vspace{0.1cm}\textbf{【有趣的小知识】}}\vspace{0.1cm}\end{minipage}\end{minipage}\end{center}\vspace{0.2cm}}{}
\newenvironment{figplaceholder}[1]{\vspace{0.2cm}\begin{center}\fbox{\begin{minipage}{0.82\textwidth}\centering\textbf{【插图位置】}\\[0.2cm]#1\end{minipage}}\end{center}\vspace{0.2cm}}{}

\title{\Huge\textbf{焊接技术小故事}}
\date{}
\begin{document}\maketitle\thispagestyle{empty}

\section{金属是怎么"粘"在一起的？}

小明看到爸爸在修家里的大门——铁门的一根栏杆断了。爸爸拿出一根焊条，用一个像大钳子的东西（焊钳）夹住焊条，在断裂处"划"了一下——啪的一声，一道刺眼的亮光闪过。几十秒后，断开的铁栏杆被"焊"上了。

小明问："为什么不用胶水？"爸爸说："胶水粘不住铁——至少不够牢固。焊接是把两块金属加热到熔化，让它们变成一整块。"

这就是焊接的本质——\textbf{冶金结合}。和胶水（物理黏附）完全不同：胶水是两个表面之间用一层不同材料粘起来；焊接是两个表面的材料本身融为一体。

\begin{figplaceholder}
插图：对比示意图——左半\"
胶水黏附"（两个金属块中间夹一层胶水）vs 右半"焊接熔合"（两个金属块在接口处融为一体）
\end{figplaceholder}

焊接是现代社会的基础——没有焊接，就没有摩天大楼（钢结构）、没有汽车（车身焊接）、没有轮船（船体分段焊接）、没有管道（油气管道焊接）。世界上使用的钢材中，约有一半需要在某处经过焊接才能变成有用产品。

\begin{thinkbox}
观察一下你周围的东西——找找哪些东西是由焊接制造的。学校的铁栅栏、大桥的钢梁、自行车的车架……试试看能不能找到焊接的痕迹（一条细长的"焊缝"）。
\end{thinkbox}

\begin{funfact}
世界上最长的焊接结构是什么？——布良科-欧亚大陆油气管道（全长约8000公里），无数段钢管被环焊连接在一起。如果把所有的焊缝加起来——总长度超过数百公里。
\end{funfact}

\newpage

\section{为什么焊工要戴面罩？}

小明看见爸爸焊接时戴着一个黑色的面罩——只露出一个深色的小窗口，什么都看不见。

"为什么需要这个？"

"因为电弧光太强了——会烧伤眼睛。"

焊接时焊条和金属之间产生\textbf{电弧}——温度高达约5000-6000°C（比太阳表面还高）。电弧除了产生热外，还会发出强烈的\textbf{紫外线和红外线}。紫外线会"烧伤"角膜（就像雪盲症——雪地反射的紫外线同样会伤害眼睛），红外线会加热和损伤视网膜。

焊工面罩的滤光片（也叫"焊接滤光镜"）有特殊的镀层——可以阻挡99\%以上的紫外线和红外线，只让少量可见光通过。

\begin{figplaceholder}
插图：焊工面罩的剖面图——标注滤光片/紫外线阻挡/红外线阻挡的不同镀层
\end{figplaceholder}

同样的道理——日食眼镜也是用类似原理（阻挡绝大部分有害光线、只保留少量可见光）。但焊接面罩的遮光号（从9号到14号）比日食眼镜更深——因为电弧比太阳（你直视它时）更亮。

\begin{thinkbox}
为什么"被弧光闪了一下眼睛"非常危险？——电弧的紫外线和电焊时溅起的熔融金属星点都可能造成永久伤害。焊工被弧光"闪了眼"后6-12小时才会开始疼痛（角膜上皮细胞的紫外线损伤需要时间显现症状）。所以永远不要在没有防护的情况下看焊接！
\end{thinkbox}

\begin{funfact}
2009年，一个荷兰艺术家用焊接面罩的滤光片做了一个"装置艺术"——透过它看太阳，可以看到完美的日食倒影。这个作品叫"The Solar Eclipse in a Welder's Helmet"。
\end{funfact}

\newpage

\section{太空里能焊接吗？}

在太空中焊接——听起来很科幻，但实际上1969年苏联宇航员就在太空中做过焊接实验了。为什么要在太空里焊接？

因为在太空里维修航天器、组装空间站、修复卫星——都需要焊接。螺栓连接在频繁的冷热交替（轨道上每90分钟经历一次-150°C到+150°C的循环）中可能松动——焊接则不会。

\textbf{太空焊接的难点：}
\begin{enumerate}
  \item \textbf{没有重力：}熔化后的金属不会"往下流"——它会变成悬浮的小球，四处飘散
  \item \textbf{真空：}地球上的气体保护焊（MIG/TIG）在太空中不能用——保护气体会扩散到真空中去
  \item \textbf{温度极端：}向阳面的高温和背阴面的低温同时存在——焊接区的热应力非常复杂
\end{enumerate}

解决办法：用\textbf{电子束焊接（EBW）}——在真空中用高速电子束（不是电弧）加热金属。电子束聚焦精准、能量集中，而且不需要保护气体——恰恰适合太空环境。

\begin{figplaceholder}
插图：电子束焊接的原理——电子枪发射高速电子→电磁透镜聚焦→电子束撞击工件→局部熔化
\end{figplaceholder}

未来的太空制造：国际空间站上已经有一台3D打印机（塑料的），但金属3D打印和焊接还需要更多技术成熟度。当人类开始建造月球基地和火星飞船时——太空焊接将是不可或缺的技术。

\begin{thinkbox}
如果你要在太空里修理一个损坏的卫星——你选择焊接还是机械连接（螺栓/卡扣）？各自的优缺点是什么？
\end{thinkbox}

\begin{funfact}
2024年，欧洲航天局（ESA）宣布了一个项目——在太空中用"摩擦搅拌焊"（FSW）来连接铝合金板材，用于制造未来的大型空间站结构。摩擦搅拌焊不需要熔化金属就可以焊接——它通过高速旋转的搅拌头产生摩擦热，让材料在固态下"揉"在一起。没有"熔化→凝固"的过程，所以没有气孔、没有裂纹、变形很小。
\end{funfact}

\end{document}
